\SourceLeaf{262}\Article{7}
Durch Benutzung der Punkte $E,E^*$ lässt sich die Abhängigkeit des letzten Weges des Lichtstrahls von dem ersten einfach so ausdrücken: der letzte Weg hat gegen den Punkt $E^*$ dieselbe Lage, welche der nur einmal gebrochene Lichtstrahl gegen $E$ haben würde, wenn in $E$ sich eine brechende Fläche mit dem Halbmesser $\frac{n^0-n^*}{k}$ befände, durch welche der Lichtstrahl aus dem ersten Mittel unmittelbar in das letzte Mittel überginge. Dies gilt für den Fall, wo das erste und das letzte Mittel ungleich sind. Sind sie hingegen gleich, oder $n^*=n^0$, wie bei Brechung durch ein oder mehrere Linsengläser, so hat der letzte Weg gegen $E^*$ dieselbe Lage, welche er gegen $E$ vermöge der Brechung durch eine in $E$ befindliche unendlich dünne Linse von der Brennweite $-\frac{n^0}k$ haben würde. Mit andern Worten: es ist verstattet, anstatt des Überganges aus dem ersten Mittel in das letzte vermöge mehrerer Brechungen, den Übergang entweder durch eine einzige Brechung, oder durch eine einzige Linse von unendlich kleiner Dicke zu substituiren, je nachdem das erste und das letzte Mittel ungleich oder gleich sind, indem man im ersten Fall der brechenden Fläche den Halbmesser $\frac{n^0-n^*}{k}$, im zweiten der Linse die Brennweite $-\frac{n^0}k$ gibt, die brechende Fläche oder die Linse in $E$ annimmt, und in beiden Fällen die Lage des ausfahrenden Strahls so viel verschiebt, als die Entfernung des Punktes $E^*$ von $E$ beträgt. Das Zeichen des Halbmessers der brechenden Fläche ist übrigens so zu verstehen, wie oben Art. 1, und das Zeichen der Brennweite so, wie weiter unten Art. 9 bemerkt werden wird.

Wegen dieser Bedeutsamkeit der Punkte $E,E^*$ scheinen diese eine besondere Benennung wohl zu verdienen: ich werde sie die \textit{Hauptpunkte} des Systems von Mitteln, oder der Linse, oder des Systems von Linsen, worauf sie sich beziehen, nennen; $E$ den ersten, $E^*$ den zweiten Hauptpunkt. Unter Ebenen der Hauptpunkte werden die durch dieselben normal gegen die Axe der $x$ gelegten Ebenen verstanden werden,

\Article{8}
Rücksichtlich der Punkte $F,F^*$ zeigen die Formeln des 6. Artikels, dass allen einfallenden Lichtstrahlen, die durch den Punkt $F$ gehen, ausfahrende entsprechen, die mit der Axe parallel sind; einfallenden hingegen, die mit der Axe parallel sind, solche ausfahrende, die sich in dem Punkte $F^*$ kreuzen; für
\SourceLeaf{263}Strahlen, die von der entgegengesetzten Seite herkommen, vertauschen diese Punkte ihre Functionen. Wenn wir also dem für einzelne Linsen bestehenden Sprachgebrauche eine erweiterte Ausdehnung geben, so können $F,F^*$ die Brennpunkte des Systems von Mitteln oder von Linsen, worauf sie sich beziehen, genannt werden, $F$ der erste, $F^*$ der zweite; die durch diese Punkte normal gegen die Axe der $x$ gelegten Ebenen mögen die Brennpunktsebenen heissen. Jene Formeln des 6. Art. zeigen zugleich, dass allen Strahlen, die sich in irgend einem andern Punkte der ersten Brennpunktsebene kreuzen, ausfahrende entsprechen, die gegen die Axe geneigt, aber unter sich parallel sind, und umgekehrt, dass allen unter sich aber nicht mit der Axe parallelen einfallenden Strahlen solche ausfahrende entsprechen, die sich in einem von $F^*$ verschiedenen Punkte der zweiten Brennpunktsebene kreuzen.

\Article{9}
Mit Hülfe dieser vier Ebenen gelangen wir zu einer sehr einfachen Construction für die Lage des ausfahrenden Strahls.

Es schneide der einfallende Strahl die erste Brennpunktsebene in dem Punkte $(1)$, die erste Hauptebene in dem Punkte $(2)$; eine Parallele mit $(1)(2)$ durch $F$ gezogen treffe die erste Hauptebene in $(3)$; eine Parallele mit der Axe durch $(2)$ treffe die zweite Hauptebene in $(4)$; endlich eine Parallele mit der Axe durch $(3)$ treffe die zweite Brennpunktsebene in $(5)$. Dann gibt $(4)(5)$ oder $(5)(4)$ die Lage des ausfahrenden Strahls. Es sind nemlich die Werthe der Coordinaten
\[
\renewcommand{\arraystretch}{1.25}
\begin{array}{c||c|c|c}
\text{für}&x&y&z\\\hline
(1)&F&B'&C'\\
(2)&E&B&C\\
F&F&0&0\\
(3)&E&B-B'&C-C'\\
(4)&E^*&B&C\\
(5)&F^*&B-B'&C-C'
\end{array}
\]
Aus den Formeln des 6. Art. folgt also, dass der ausfahrende Strahl durch $(4)$ und $(5)$ geht; das erstere unmittelbar, das andere, weil
\SourceLeaf{264}
\[
\begin{aligned}
B-B'&=\frac{\beta^0}{n^0}(E-F)=-\frac{\beta^0}k\\
C-C'&=\frac{\gamma^0}{n^0}(E-F)=-\frac{\gamma^0}k
\end{aligned}
\]

In dem gewöhnlichsten Falle, wo $n^*=n^0$, also $F^*-E^*=E-F$, wird die Construction noch einfacher, weil der Punkt $(3)$ überflüssig wird; man braucht nur $(1),(2),(4)$ wie vorhin zu bestimmen, und $(4)(5)$ mit $(1)E$ parallel zu ziehen.

Geht die Richtung des einfallenden Strahls durch $E$, so geht allemal die Richtung des ausfahrenden durch $E^*$, und ist, in dem Falle, wo $n^*=n^0$ ist, zugleich jener parallel. Man pflegt (bei einfachen Linsen) einen solchen Strahl einen Hauptstrahl zu nennen.

Die Entfernungen der zweiten Brennpunktsebene von der zweiten Haupt\-ebene, und der ersten Hauptebene von der Ebene des ersten Brennpunkts, oder die Grössen $-\frac{n^*}k,-\frac{n^0}k$ könnte man die \textit{Brennweiten} des Systems der Mittel nennen, wenn es nicht angemessener schiene, den Gebrauch dieser Benennung auf den Fall zu beschränken, wo das letzte Mittel dasselbe ist, wie das erste, also jene Entfernungen unter sich gleich sind. Um dem gewöhnlichen Sprachgebrauche conform zu bleiben, sehen wir die Brennweite als positiv an, wenn dem ersten Hauptpunkte eine grössere Coordinate entspricht, als dem ersten Brennpunkte, so, dass die Brennweite immer durch $-\frac{n^0}k=-\frac{n^*}k$ ausgedrückt wird.

\Article{10}
In den oben Art. 4 für den Platz des Bildes gegebenen Formeln ist es, wie man leicht sieht, verstattet, anstatt $N^0,N^*$ andere Punkte zu setzen, wenn man nur zugleich anstatt $g,h,k,l$ die entsprechenden $G,H,K,L$ substituirt. Indem wir dazu die Hauptpunkte wählen, erhalten wir folgende Ausdrücke:
\[
\begin{aligned}
\xi^*&=E^*-\frac{n^*(E-\xi)}{n^0+k(E-\xi)}\\
\eta^*&=\frac{n^0\eta}{n^0+k(E-\xi)}\\
\zeta^*&=\frac{n^0\zeta}{n^0+k(E-\xi)}
\end{aligned}
\]
Der ersten Formel kann man auch folgende Gestalt geben
\[
\frac{n^*}{\xi^*-E^*}+\frac{n^0}{E-\xi}=-k
\]
\SourceLeaf{265}Wählen wir die Brennpunkte, so erhalten wir
\[
\begin{aligned}
\xi^*&=F^*+\frac{n^0n^*}{kk(F-\xi)}\\
\eta^*&=\frac{n^0\eta}{k(F-\xi)}\\
\zeta^*&=\frac{n^0\zeta}{k(F-\xi)}
\end{aligned}
\]

Wegen des häufigen Gebrauchs mögen die Formeln auch noch in der Gestalt hier stehen, die sie annehmen, wenn das erste und das letzte Mittel gleich sind, und die Brennweite mit $\varphi$ bezeichnet wird.
\[
\begin{gathered}
\frac1{\xi^*-E^*}+\frac1{E-\xi}=\frac1\varphi\\
(\xi^*-F^*)(F-\xi)=\varphi\varphi\\
\eta^*=-\frac{\varphi\eta}{F-\xi}=-\frac{\eta(\xi^*-F^*)}{\varphi}\\
\zeta^*=-\frac{\varphi\zeta}{F-\xi}=-\frac{\zeta(\xi^*-F^*)}{\varphi}
\end{gathered}
\]

\Article{11}
Die vier Hülfspunkte $E,E^*,F,F^*$ verlieren ihre Anwendbarkeit in dem besondern Falle, wo $k=0$ ist, also jene Punkte als unendlich entfernt von den brechenden Flächen betrachtet werden müssten. Man kann sich in diesem Falle unmittelbar an die allgemeinen zur Auflösung der Hauptaufgaben oben mitgetheilten Formeln halten, welche hier folgende Gestalt annehmen.

Wenn die Gleichungen für den einfallenden Strahl so ausgedrückt sind
\[
\begin{aligned}
y&=\frac{\beta^0}{n^0}(x-N^0)+b^0\\
z&=\frac{\gamma^0}{n^0}(x-N^0)+c^0
\end{aligned}
\]
so sind die für den ausfahrenden
\[
\begin{aligned}
y&=\frac{l\beta^0}{n^*}.(x-N^*)+gb^0+h\beta^0\\
z&=\frac{l\gamma^0}{n^*}.(x-N^*)+gc^0+h\gamma^0
\end{aligned}
\]
Setzt man zur Abkürzung
\[
N^*-\frac{hn^*}{l}=N^{**}
\]
oder, was dasselbe ist, weil $gl=1$,
\SourceLeaf{266}
\[
N^*-ghn^*=N^{**}
\]
so erscheinen diese Formeln noch einfacher, nemlich
\[
\begin{aligned}
y&=\frac{l\beta^0}{n^*}.(x-N^{**})+gb^0\\
z&=\frac{l\gamma^0}{n^*}.(x-N^{**})+gc^0
\end{aligned}
\]
Für den Platz des Bildes desjenigen Punktes, dessen Coordinaten $\xi,\eta,\zeta$ sind, erhalten wir die Coordinaten
\[
\begin{aligned}
\xi^*&=N^*-ghn^*-\frac{n^*}{n^0}gg(N^0-\xi)\\
&=N^{**}-\frac{n^*}{n^0}gg(N^0-\xi)\\
\eta^*&=g\eta\\
\zeta^*&=g\zeta
\end{aligned}
\]
Es erhellt hieraus, dass der Punkt der Axe der $x$, welcher in Gemässheit der von uns immer gebrauchten Bezeichnungsart mit $N^{**}$ zu bezeichnen ist, das Bild des Punktes $N^0$ vorstellt, und dass das Linearverhältniss der Theile eines zusammengesetzten Bildes zum Object constant, nemlich wie $g$ zu $1$ oder wie $1$ zu $l$ ist.

\Article{12}
Der im vorhergehenden Artikel betrachtete Fall kommt vor bei einem Fernrohre, dessen Gläser für ein weitsichtiges Auge und für das deutliche Sehen unendlich entfernter Gegenstände gestellt sind. Aus obigen Formeln erhellt, dass die Richtung des ausfahrenden Strahls bloss von der Richtung des einfallenden abhängt, dass also parallel unter sich einfallenden Strahlen auch parallel ausfahrende entsprechen, und dass die Tangente der Neigung der erstern gegen die Axe sich zu der Tangente der Neigung der letztern verhält, wie $1$ zu $l$. Die Zahl $l=\frac1g$ ist also das, was man die Vergrösserung des Fernrohrs nennt, und ihr positives oder negatives Zeichen bedeutet die aufrechte oder verkehrte Erscheinung. Lässt man die einfallenden und ausfahrenden Strahlen ihre Funktionen vertauschen, indem man den Gegenständen die Ocularseite zuwendet. so erscheinen sie in demselben Verhältnisse verkleinert, und hierauf gründet sich das eben so bequeme als scharfe Verfahren zur Bestimmung der Vergrösserung eines Fern\SourceLeaf{267}rohrs, welches ich 1823 im 2. Bande der \textit{Astronomischen Nachrichten} mitgetheilt habe.

Eine andere Methode, die Vergrösserung zu bestimmen, beruht auf der Vergleichung eines Gegenstandes mit seinem Bilde nach dem linearen Verhältnisse. \textsc{Ramsdens} Dynameter ist nichts anderes, als eine Vorrichtung, den Durchmesser des in $N^{**}$ fallenden Bildes von der kreisrunden Begrenzung des Objectivs zu messen, wobei man sich natürlich erst vergewissern muss, dass dieses Bild wirklich erscheint und nicht etwa durch eine innere Blendung verdeckt ist. Auch muss das Bild ein reelles sein, wozu erforderlich ist, dass $ghn^*$ negativ wird: bei einem \textsc{Galileis}chen Fernrohr, wo dieses Bild nur ein virtuelles ist, würde man ein genaues Resultat nur mit einem mikrometrischen Mikroskope erlangen können, welches auch in \textit{allen} Fällen. wo man eine grössere Schärfe wünscht, den Vorzug verdienen würde. Übrigens erhellt aus dem vorhergehenden Artikel, dass eben so gut ein schickliches vom Objectiv entferntes Object gebraucht werden kann, so lange nur die Entfernung nicht so gross wird, dass das Bild aufhört ein reelles oder mit dem Mikroskope erreichbares zu sein. Endlich mag noch bemerkt werden, dass der Punkt $N^{**}$ derjenige ist, welcher in der Theorie der Fernröhre mit der Benennung \textit{Ort des Auges} belegt wird.

\Article{13}
Um die allgemeinen Vorschriften des 2. Artikels auf den Fall einer einfachen Glaslinse anzuwenden, bezeichnen wir das Brechungsverhältniss beim Übergange aus Luft in Glas mit $n:1$; die Halbmesser der ersten und zweiten Fläche mit $(n-1)f$ und $(n-1)f'$; die Dicke der Linse mit $ne$. Wir haben also anstatt der dortigen Bezeichnungen
\[
\begin{array}{ccc}
n^0&\mathit{hier}&1\\
n'&.\,.\,.\,.&n\\
n''&.\,.\,.\,.&1\\
t'&.\,.\,.\,.&e\\
u^0&.\,.\,.\,.&-\dfrac1f\\
u'&.\,.\,.\,.&-\dfrac1{f'}
\end{array}
\]
und folglich
\SourceLeaf{268}
\[
\begin{aligned}
g&=1+u^0t'=\frac{f-e}{f}\\
h&=t'=e\\
k&=u^0+u'+t'u^0u'=-\frac{f+f'-e}{ff'}\\
l&=1+u't'=\frac{f'-e}{f'}
\end{aligned}
\]
Für die Brennweite $\varphi$ haben wir also nach Art. 9
\[
\varphi=\frac{ff'}{f+f'-e}
\]
für die beiden hier mit $E,E'$ zu bezeichnenden Hauptpunkte nach Art. 6
\[
\begin{aligned}
E&=N^0+\frac{ef}{f+f'-e}=N^0+\frac{e\varphi}{f'}\\
E'&=N'-\frac{ef'}{f+f'-e}=N'-\frac{e\varphi}{f}
\end{aligned}
\]
und für die beiden Brennpunkte $F,F'$
\[
\begin{aligned}
F&=E-\varphi=N^0-\frac{f(f'-e)}{f+f'-e}\\
F'&=E'+\varphi=N'+\frac{f'(f-e)}{f+f'-e}
\end{aligned}
\]
Für den Durchschnittspunkt der (nöthigenfalls vorwärts oder rückwärts verlängerten) geraden Linie, welche ein Hauptstrahl im Innern der Linse beschreibt, mit der Axe findet man leicht
\[
x=N^0+\frac{nef}{f+f'}=N'-\frac{nef'}{f+f'}
\]
Diesen Punkt, welcher also von der Neigung des Hauptstrahls unabhängig ist, nennen einige Schriftsteller den optischen Mittelpunkt der Linse, eine Auszeichnung, welche dieser sonst gar keine merkwürdigen Eigenschaften darbietende Punkt kaum verdient haben möchte, und die hie und da zu dem Irrthum verleitet zu haben scheint, als ob die einfachen Relationen zwischen Bild und Object, welche bei einer unendlich dünnen Linse Statt finden, sich auf eine Linse von endlicher Dicke bloss durch Beziehung auf jenen Mittelpunkt übertragen liessen, während diese Übertragung, wie oben gezeigt ist, nur dann gültig ist, wenn das Object auf den ersten, das Bild auf den zweiten Hauptpunkt bezogen wird. Bei einem Systeme von mehrern Linsen, also schon bei einem achromatischen Doppelobjective, kann ohnehin von einem Mittelpunkte in jenem Sinne gar nicht die Rede sein. Will man die Benennung beibehalten, so würde ich für angemessener halten, sie demjenigen Punkte beizulegen, welcher zwischen den beiden
\SourceLeaf{269}Hauptpunkten (mithin auch zwischen den beiden Brennpunkten) in der Mitte liegt, und der mit jenem Punkte nur dann zusammenfällt, wenn die Linse gleichseitig ist. Dieser Punkt hat die praktisch nützliche Eigenschaft, durch Umwenden der Linse leicht und mit Schärfe bestimmbar zu sein; denn offenbar ist es dieser Punkt, der beim Umwenden wieder den vorigen Platz einnehmen muss, wenn der Platz des Bildes von einem festen Objecte ungeändert bleiben soll.

Es mag hier noch bemerkt werden, dass die Entfernung der beiden Hauptpunkte von einander
\[
E'-E=ne-\frac{e(f+f')}{f+f'-e}=(n-1)e-\frac{ee}{f+f'-e}
\]
wird, also, insofern gewöhnlich $e$ gegen $f+f'-e$ sehr klein ist, von $(n-1)e$ oder von der durch $\frac{n-1}{n}$ multiplicirten Dicke der Linse kaum merklich verschieden ist.

\Article{14}
An die Stelle der allgemeinen Formeln des 2. Art., durch welche aus dem Wege des einfallenden Lichtstrahls der Weg des ausfahrenden bestimmt wird, lassen sich für den Fall eines Systems von Linsen auf einer gemeinschaftlichen Axe bequemere setzen, indem man, anstatt der Halbmesser der einzelnen brechenden Flächen und ihrer gegenseitigen Abstände, die Brennweiten der einzelnen Linsen und die Entfernungen ihrer zweiten Hauptpunkte von den ersten der folgenden Linsen einführt. Die neuen Formeln werden denen des 2. Art. ganz ähnlich, enthalten aber nur halb so viele Elemente. Da ihre Ableitung aus dem Vorhergehenden sehr leicht ist, so wird es hinreichend sein, sie in gebrauchfertiger Form hieher zu setzen.

Wir bezeichnen die Brennweiten der einzeln auf einander folgenden Linsen mit $\varphi^0,\varphi',\varphi''$ u. s. f.; ihre Hauptpunkte hier, abweichend von der bisherigen Bezeichnungsart, die ersten mit $E^0,E',E''$ u. s. f., die zweiten mit $I^0,I',I''$ u. s. f. Zur Abkürzung schreiben wir
\[
\begin{aligned}
-\frac1{\varphi^0}&=u^0,&-\frac1{\varphi'}&=u',&-\frac1{\varphi''}&=u''\quad\text{u. s. f.}\\
E'-I^0&=t',&E''-I'&=t'',&E'''-I''&=t'''\quad\text{u. s. f.}
\end{aligned}
\]
die letzten Glieder in diesen Reihen mögen als solche durch ein Sternchen ausgezeichnet werden.

\SourceLeaf{270}Setzt man nun die Gleichungen für den einfallenden Strahl in die Form
\[
\begin{aligned}
y&=\beta^0(x-E^0)+b^0\\
z&=\gamma^0(x-E^0)+c^0
\end{aligned}
\]
für den ausfahrenden hingegen in folgende
\[
\begin{aligned}
y&=\beta^*(x-I^*)+b^*\\
z&=\gamma^*(x-I^*)+c^*
\end{aligned}
\]
so wird, wenn die vier Grössen $g,h,k,l$ durch Formeln bestimmt werden, die mit den im 2. Art. als (5) bezeichneten ganz identisch sind,
\[
\begin{aligned}
b^*&=gb^0+h\beta^0,&c^*&=gc^0+h\gamma^0\\
\beta^*&=kb^0+l\beta^0,&\gamma^*&=kc^0+l\gamma^0
\end{aligned}
\]
Für die beiden Hauptpunkte des Linsensystems, als Ganzes betrachtet, wird
\[
\begin{aligned}
\text{für den ersten}\quad x&=E^0-\frac{1-l}{k}\\
\text{für den zweiten}\quad x&=I^*+\frac{1-g}{k}
\end{aligned}
\]
Ferner wird für die beiden Brennpunkte des Linsensystems
\[
\begin{aligned}
\text{für den ersten}\quad x&=E^0+\frac{l}{k}\\
\text{für den zweiten}\quad x&=I^*-\frac{g}{k}
\end{aligned}
\]
\[
\text{die Brennweite selbst ist}\quad=-\frac1k
\]

Die Formeln für den Fall, wo das System nur aus zwei Linsen besteht, verdienen noch besonders hergeschrieben zu werden. Man hat nemlich
\[
\begin{aligned}
g&=\frac{\varphi^0-t'}{\varphi^0}\\
h&=t'\\
k&=-\frac{\varphi^0+\varphi'-t'}{\varphi^0\varphi'}\\
l&=\frac{\varphi'-t'}{\varphi'}
\end{aligned}
\]
Die Werthe von $x$ für die beiden Hauptpunkte sind
\SourceLeaf{271}
\[
E^0+\frac{t'\varphi^0}{\varphi^0+\varphi'-t'}\quad\text{und}\quad I'-\frac{t'\varphi'}{\varphi^0+\varphi'-t'}
\]
und die Brennweite
\[
=\frac{\varphi^0\varphi'}{\varphi^0+\varphi'-t'}
\]
Man sieht, dass diese Formeln denen ganz analog sind, die im 13. Artikel für die Bestimmung der Hauptpunkte und der Brennweite einer einfachen Linse gegeben sind, indem an die Stelle der dortigen $f^0,f',e$ hier die Grössen $\varphi^0,\varphi',t'$ treten.

Die Entfernung der beiden Hauptpunkte von einander wird in dem Fall zweier Linsen
\[
\begin{aligned}
&=I'-E^0-\frac{t'(\varphi^0+\varphi')}{\varphi^0+\varphi'-t'}\\
&=I^0-E^0+I'-E'-\frac{t't'}{\varphi^0+\varphi'-t'}
\end{aligned}
\]
Ist $t'$ sehr klein, wie bei achromatischen Doppellinsen von der gewöhnlichen Einrichtung immer der Fall ist, so wird das letzte Glied unbedeutend, und daher die Entfernung der beiden Hauptpunkte von einander für eine solche Doppellinse als Ganzes betrachtet sehr nahe der Summe der beiden Werthe gleich, welche diese Entfernung in den Linsen, einzeln genommen, hat.

Übrigens ist von selbst klar, dass die sämmtlichen in dem gegenwärtigen Artikel aufgeführten Formeln ohne alle Veränderung auf den Fall übertragen werden können, wo anstatt einfacher Linsen partielle Systeme von Linsen zu Einem ganzen Systeme vereinigt werden sollen.

\Article{15}
Die optischen Erscheinungen sowohl durch eine einfache Linse, als durch ein System von mehreren auf gemeinschaftlicher Axe, hängen, wie wir gezeigt haben, von drei Elementen ab, welche durch das Brechungsverhältniss (oder durch die Brechungsverhältnisse, wenn sie für die verschiedenen Linsen verschieden sind), und die Lagen und Halbmesser der brechenden Flächen bestimmt sind: da jedoch diese Grössen gewöhnlich unmittelbar nicht bekannt sind, so bleibt noch übrig, einiges über die Methode zu sagen, durch welche umgekehrt aus beobachteten Erscheinungen jene drei Elemente abgeleitet werden können. Wir bezeichnen die verschiedenen hiebei in Frage kommenden Punkte der Axe auf folgende Weise:
